\subsection{Jacobson rings}

DEFINITION 1. A ring \(\mathbf{A}\) is called a Jacobson ring if every prime ideal of \(\mathbf{A}\) is the intersection of a family of maximal ideals.

\textbf{Examples}\par

\begin{enumerate}
\def\labelenumi{(\arabic{enumi})}
\item
  Every field is a Jacobson ring.
\item
  The ring Z is a Jacobson ring, the unique prime ideal which is not maximal (0) being the intersection of the maximal ideals \((p)\) of Z, where p runs through the set of prime numbers (cf.~Proposition 4).
\item
  Let A be a Jacobson ring and let a be an ideal of A. Then A/a is a Jacobson ring, for the ideals of A/a are of the form b/a, where b is an ideal of A containing a and b/a is prime (resp. maximal) if and only if b is.
\end{enumerate}

PROPOSITION 3. For a ring A to be a Jacobson ring, it is necessary and sufficient that, for every ideal a of A, the Jacobson radical of A/a be equal to its nilradical (Chapter II, § 2, no. 6).

The Jacobson radical (resp. nilradical) of A/a is the intersection of the maximal (resp. prime) ideals of A/a (Algebra, Chapter VIII, § 6, no. 3, Definition 3 and Commutative Algebra, Chapter II, § 2, no. 6, Proposition 13). The stated condition means therefore that for every ideal a of A the intersection of the prime ideals containing a is equal to the intersection of the maximal ideals containing a. This condition obviously holds for every ideal a of A if A is a Jacobson ring; conversely, if it holds for every prime ideal of A, A is a Jacobson ring by definition.

COROLLARY. Let A be a Jacobson ring; for every ideal a of A, the radical of a is the intersection of the maximal ideals of A containing a.

It is sufficient to note that A/a is a Jacobson ring.

PROPOSITION 4. Let A be a principal ideal domain and \((p_{\lambda})_{\lambda \in L}\) a representative system of extremal elements of A (Algebra, Chapter VII, § 1, no. 3, Definition 2). For A to be a Jacobson ring, it is necessary and sufficient that L be infinite.

The maximal ideals of A are the \(Ap_{\lambda}\) (loc. cit., no. 2, Proposition 2). If L is finite, their intersection is the ideal Ax, where \(x = \prod_{\lambda \in L} p_{\lambda}\) (ibid.) and hence different from (0); on the other hand, if L is infinite, the intersection of the \(Ap_{\lambda}\) is (0), every element \(\neq 0\) of A being divisible by only a finite number of extremal elements (loc. cit., no. 3, Theorem 2). The proposition then follows from the fact that (0) is the only prime ideal which is not maximal in A (Algebra, Chapter VI, §1, no. 13, Proposition 14 (DIV)).

PROPOSITION 5. Let A be a ring and B an A-algebra integral over A. If A is a Jacobson ring, so is B.

Replacing A by its canonical image in B, we may assume that \(A \subset B\) . Let \(p'\) be a prime ideal of B and let \(p = A \cap p'\) . There exists by hypothesis a family \((m_{\lambda})_{\lambda \in L}\) of maximal ideals of A whose intersection is equal to p.~For all \(\lambda \in L\) there exists a maximal ideal \(m_{\lambda}'\) of B lying above \(m_{\lambda}\) and containing \(p'\) (§ 2, no. 1, Proposition 1 and Corollary 2 to Theorem 1). If we write \(q' = \bigcap_{\lambda \in L} m_{\lambda}'\) , then \(q' \cap A = \bigcap_{\lambda \in L} m_{\lambda} = p\) and \(q' \supset p'\) , whence \(q' = p'\) (§ 2, no. 1, Corollary 1 to Proposition 1).

THEOREM 3. Let A be a Jacobson ring, B a finitely generated A-algebra and \(\rho: \mathbf{A} \to \mathbf{B}\) the canonical homomorphism. Then:

\begin{enumerate}
\def\labelenumi{(\roman{enumi})}
\item
  B is a Jacobson ring.
\item
  For every maximal ideal \(\mathfrak{m}'\) of \(\mathbf{B}\) , \(\mathfrak{m} = \rho^{-1} (\mathfrak{m}')\) is a maximal ideal of \(\mathbf{A}\) and \(\mathbf{B} / \mathfrak{m}'\) is a finite algebraic extension of \(\mathbf{A} / \mathfrak{m}\) .
\end{enumerate}

Let \(\mathfrak{p}'\) be a prime ideal of B and \(\mathfrak{p} = \rho^{-1} (\mathfrak{p}')\) . Let \(v\) be an element \(\neq 0\) of \(\mathrm{B / p'}\) .

As \(B/p'\) is a finitely generated integral \((A/p)\) -algebra and the canonical homomorphism \(\phi: A/p \to B/p'\) is injective, there exists an element \(u \neq 0\) of \(A/p\) such that, for every homomorphism \(f\) from \(A/p\) to an algebraically closed field \(L\) whose kernel does not contain \(u\) , there exists a homomorphism \(g\) from \(B/p'\) to \(L\) whose kernel does not contain \(v\) and for which \(f = g \circ \phi\) (no. 1, Corollary 3 to Theorem 1). Since \(A\) is a Jacobson ring, there exists a maximal ideal \(m\) of \(A\) containing \(p\) and such that \(u \notin m/p\) . We take \(L\) to be an algebraic closure of \(A/m\) and \(f\) to be the canonical homomorphism \(A/p \to L\) ; let

\[
g \colon \mathrm{B} / \mathfrak {p} ^ {\prime} \rightarrow \mathrm{L}
\]

be a homomorphism such that \(f = g \circ \phi\) and \(g(v) \neq 0\) . Then

\[
\mathbf {A} / \mathfrak {m} \subset g (\mathbf {B} / \mathfrak {p} ^ {\prime}) \subset \mathbf {L},
\]

hence \(g(\mathbf{B} / \mathfrak{p}')\) is a subfield of \(\mathbf{K}\) (Algebra, Chapter V, § 3, no. 2, Proposition 3) and the kernel of \(g\) is therefore a maximal ideal of \(\mathbf{B} / \mathfrak{p}'\) not containing \(v\) . Thus it is seen that the intersection of the maximal ideals of \(\mathbf{B} / \mathfrak{p}'\) is reduced to 0, which proves that \(\mathbf{B}\) is a Jacobson ring. Moreover, if \(\mathfrak{p}'\) is maximal, \(g\) is necessarily injective and hence \(\mathfrak{p} = \mathfrak{m}\) is maximal; finally \(\mathbf{B} / \mathfrak{p}'\) is then a finitely generated algebra over the field \(\mathbf{A} / \mathfrak{m}\) and hence is a finite extension of \(\mathbf{A} / \mathfrak{m}\) (no. 1, Corollary 2 to Theorem 1).

COROLLARY 1. Every finitely generated algebra A over Z is a Jacobson ring; for a prime ideal p of A to be maximal, it is necessary and sufficient that the ring A/p be finite.

If the integral domain A/p is finite, it is a field, as, for every \(u \neq 0\) in A/p, the mapping \(v \mapsto uv\) of A/p to itself is injective and hence bijective since A/p is finite. Conversely, for every maximal ideal m of A, the inverse image of m in Z is a maximal ideal (p) and A/m is finite over the prime field \(\mathbf{Z}/(\mathfrak{p}) = \mathbf{F}_{\mathfrak{p}}\) by Theorem 3.

COROLLARY 2. Let \((\mathbf{P}_{\lambda})_{\lambda \in \mathbf{L}}\) be a family of polynomials in \(\mathbf{Z}[\mathbf{X}_1,\ldots ,\mathbf{X}_n]\) and let \(\mathbf{Q}\) be a polynomial in \(\mathbf{Z}[\mathbf{X}_1,\dots ,\mathbf{X}_n]\) such that, for every system of elements \((x_i)_{1\leqslant i\leqslant n}\) belonging to a finite field and for which \(\mathrm{P}_{\lambda}(x_1,\ldots ,x_n) = 0\) for all \(\lambda\) , also \(\mathrm{Q}(x_1,\ldots ,x_n) = 0\) . Then, if \(\mathfrak{a}\) is the ideal of \(\mathbf{Z}[\mathbf{X}_1,\ldots ,\mathbf{X}_n]\) generated by the \(\mathrm{P}_{\lambda}\) , there exists an integer \(m > 0\) such that \(\mathbf{Q}^m\in \mathfrak{a}\) . Moreover, for every reduced ring R and every system \((y_i)_{1\leqslant i\leqslant n}\) of elements of R such that \(\mathrm{P}_{\lambda}(y_1,\ldots ,y_n) = 0\) for all \(\lambda\) , also \(\mathrm{Q}(y_1,\ldots ,y_n) = 0\) .

The second assertion follows from the first since the ideal of \(\mathbf{Z}[\mathrm{X}_1,\ldots ,\mathrm{X}_n]\) consisting of the polynomials \(\mathbf{P}\) such that \(\mathrm{P}(y_1,\dots,y_n) = 0\) contains a. To show the first assertion, it suffices to note that, for every maximal ideal m of \(\mathbf{A} = \mathbf{Z}[\mathrm{X}_1,\dots,\mathrm{X}_n]\) containing a, A/m is a finite field (Corollary 1) and the hypothesis implies that the canonical image of Q in A/m is zero; then Q belongs to the intersection of the maximal ideals of A containing a, which is the radical of a (Corollary to Proposition 3).

COROLLARY 3. Let A be a Jacobson ring. If there exists a finitely generated A-algebra B containing A and which is a field, then A is a field and B is an algebraic extension of A.

It suffices to apply Theorem 3 (ii) with \(\mathfrak{m}' = (0)\) .


